\section{Conclusion}
\label{sec:conclusion}

We presented \sys, a hardware attribution pipeline that bridges the gap between PyTorch's
operator-level abstractions and NVIDIA's kernel-level hardware performance counters.  The
core contribution is a mapping from CUDA kernels to \texttt{aten::} operators via NVTX
temporal enclosure with per-stream interval trees, invocation-order matching across
incompatible clock domains, and Inductor debug-artifact enrichment for compiled models.
The result is \texttt{profile.json}: a hardware-attributed operator profile mapping each
\texttt{aten::} operator to its constituent kernels and their full hardware counter state,
produced automatically without manual kernel-to-operator correlation.  A curated 20-counter
metric set with duration-weighted aggregation and architecture-specific fallbacks covers all
hardware bottleneck axes; GPU clock locking ensures reproducible per-operator duration
measurements; and layer deduplication reduces \ncu replay time by
$O(\text{num\_duplicate\_layers})$ for models with repeated structure.

To demonstrate that attributed profiles are actionable, we implemented a profile-guided
FX pass framework (\secref{sec:pgo}) and applied it across three case studies on an NVIDIA
RTX PRO 6000 Blackwell Server Edition.  Each optimization decision was grounded in a
specific counter reading from the attributed profile, and each claimed improvement was
verified by re-profiling at identical locked clocks:
\textbf{1.76$\times$} on GPT-2 (BF16 GEMM promotion, \texttt{tensor\_core\_active\_pct}:
0\%$\to$40.5\%);
\textbf{2.24$\times$} on SDPA Attention (QKV fusion raises occupancy 16.6\%$\to$73.9\%;
SDPA kernel 3.0$\times$); and
\textbf{3.63$\times$} on LSTM Sequence Encoder (structural re-routing from Inductor's
decomposed SIMT path to cuDNN's fused Tensor-Core RNN, diagnosable only via the attribution
data).  The LSTM case is the clearest illustration of the pipeline's value: the bottleneck
is invisible to timing-only profilers but apparent from the high unattributed fraction and
scalar GEMM signatures in the attributed profile.

\paragraph{Limitations.}
Attribution coverage degrades for cuDNN-backed operators (\texttt{nn.LSTM},
\texttt{nn.MultiheadAttention} under cuDNN) where internal kernel names carry no NVTX
correspondence and no Inductor debug metadata is available; expected unattributed rates of
80--90\% make hardware-counter-driven optimization inapplicable, and bottleneck diagnosis must
fall back to kernel-name inference.  Invocation-order matching requires a deterministic
workload: models with conditional control flow or dynamic dispatch can violate the
\texttt{(kernel\_name, invocation\_index)} invariant without producing a shape error;
\sys validates invocation counts before issuing counter assignments and raises
\texttt{InvocationCountMismatchError} on violation.  \ncu serializes kernels for replay,
making concurrent multi-stream profiles require separate per-stream runs.  All speedup
figures are point estimates from 2 measured iterations at locked clocks; run-to-run variance
is not characterized.  Optimization patterns are currently user-defined per-workload;
automatic pattern selection from counter evidence is left as future work.

\paragraph{Future work.}
Online profiling via CUPTI Perfworks streaming would eliminate the separate \ncu pass, enabling
profiling during ordinary training without a dedicated replay run.  Multi-GPU attribution with
NCCL stream synchronization events would extend the pipeline to distributed workloads.
Counter-guided neural architecture search---using per-operator hardware utilization as a fitness
signal rather than validation accuracy alone---represents a longer-term direction for hardware-aware
model design.  Finally, extending the counter name mapping and pass library to AMD ROCm (\texttt{rocprof}
+ \texttt{rocm-smi}) would generalize \sys beyond NVIDIA hardware.
